\chapter{부품 정격 표}

모델과 회로가 판정 한계 안에서 동작해도 실제 부품이 그 스트레스를 받지 못하면 설계는 성립하지 않는다. 이 장의 표는 확정 부품표가 아니라 모의 스트레스와 데이터시트 값을 대조한 조건부 후보 선정이다(검증 수준 ①′). 대조는 세 번 하고 두 번 다시 판정했으며(E34, E37–E40), 기준은 매번 실행 전에 정했다.

\section{공통 가정과 판정 규칙}

\begin{itemize}
\item \textbf{냉각 기준.} SiC·GaN은 케이스 100 °C, IGBT와 방열판 기준 열저항을 주는 SiC 모듈은 방열판 65 °C.
\item \textbf{온도 한도.} 정상 운전은 T\textsubscript{J} ≤ T\textsubscript{J,max} − 25 °C, 단발 고장 사건은 T\textsubscript{J} ≤ 데이터시트 SOA의 한계.
\item \textbf{반도체 공통.} 모델 최대 차단 전압 ≤ 0.8 × 정격(기생 인덕턴스 과전압은 모델에 없음), 소자당 첨두 ≤ 0.5 × 펄스 정격. 병렬 분담은 스위칭이면 같은 몫 × 1.1, 선형 영역(바이패스 소프트 투입)은 분담을 보장할 수 없으므로 소자 1개가 모두 받는 경우로 판정한다.
\item \textbf{열.} 정상 열저항을 쓰되 최대값이 있으면 최대값, 없으면 전형값 × 1.2로 한다(E39). 정상 열저항으로 한 사건의 평균 손실을 환산하는 선정 지표이며, 짧은 단일 사건에서는 보수 쪽이지만 반복 사건의 열 누적이나 지연 변화까지 보장하지 않는다. 클램프 스위치(P5)의 온도는 \textbf{선정식 환산온도}(정의는 12.5절·부록 A)다. EPC2305는 데이터시트 단일 펄스 열임피던스를 Foster 회로망으로 맞춰 전력 이력을 겹쳐 계산한다.
\item \textbf{규칙 L(E40).} 전형값만 있는 손실 인자(R\textsubscript{DS(on)}, 스위칭 에너지, 다이오드 전압 등)는 같은 데이터시트의 다른 온도에 최대값이 있으면 그 최대/전형 비를, 없으면 × 1.2를 곱한다. 예: C4MS008075U2 도통 × 1.30, CAB530M12BM3 R\textsubscript{DS(on)} × 1.33, FF2MR12KM1P R\textsubscript{DS(on)}·스위칭 에너지·열저항 × 1.2.
\end{itemize}

\Needspace{14\baselineskip}\section{최종 부품 정격 표}

{\footnotesize
\begin{longtable}{L{\dimexpr 0.08\linewidth-2\tabcolsep\relax}L{\dimexpr 0.16\linewidth-2\tabcolsep\relax}L{\dimexpr 0.235\linewidth-2\tabcolsep\relax}L{\dimexpr 0.2625\linewidth-2\tabcolsep\relax}L{\dimexpr 0.2625\linewidth-2\tabcolsep\relax}}
\caption{부품 정격 표(최종)}\label{t:parts}\\
\toprule
\tabhead \textbf{번호} & \textbf{부품} & \textbf{후보} & \textbf{스트레스(모의)} & \textbf{판정} \\
\midrule
\endfirsthead
\toprule
\tabhead \textbf{번호} & \textbf{부품} & \textbf{후보} & \textbf{스트레스(모의)} & \textbf{판정} \\
\midrule
\endhead
\bottomrule
\endlastfoot
P1 & Stage A 스위치 & Wolfspeed C4MS008075U2 × 2(위치당) & 차단 518 V. 소자 첨두(위상 첨두 ÷ 2 × 분담 1.1, A25R): 요구 안 과도 41.9 A, 모듈 고장 58.0 A(해석 C\textsubscript{fly} 46.875 µF; 후보 45.0·42.75 µF에서 위상 106.86 A·소자당 58.77 A, E44), 요구 밖(금지) 조건 121.7 A(한도 309.5 A). 정격 충전 가장 뜨거운 위치 소자 손실 12.3 W & \textbf{통과}: 102.4 °C(규칙 L) \\
P2 & 플라잉 커패시터(해석 설계 46.875 µF) & CastorCap CC7.5u1300 × 6(위상당 45.0 µF, 채택 검토 부품) & 위상 실효 52.5 A(공차 표본 60.9 A) & \textbf{통과, 단서}: 전류 정격 30 A\textsubscript{rms}는 10 kHz·30 °C 조건. 해석 설계 대비 −4.0 \%(f\textsubscript{r} +2.1 \%)라 제어 타이밍을 45.0 µF(25.516 kHz)에 맞추고 실제 45.0 µF와 −5 \% 42.75 µF(플랜트 공진 26.178 kHz)로 다시 판정해 요구 안 과도 12/12·모듈 고장 16/16 통과(76.1·106.9 A, E44). +5 \%와 전체 온도 범위는 미검증 \\
P3 & 공진 인덕터 864.6 nH & 공심 코일 지름 60 mm·3.6회·Litz 15.2 mm\textsuperscript{2} & 첨두(A25R): 요구 안 과도 76.1 A, 모듈 고장 105.5 A(해석 46.875 µF; 후보 45.0·42.75 µF에서 106.86 A), 요구 밖(금지) 조건 221.2 A(공심, 포화 없음). 실효 60.9 A & \textbf{요구·설계 추정}: 손실 3.1–4.2 W, 포화 없음 \\
P4 & V\textsubscript{A} 클램프 저항 & Ohmite 891SP 2.4 Ω × 3(R2-W460 2.7 Ω × 3) & 가지 에너지 ≤ 3,405·3,729 J(공칭 문턱), 구현 요구 범위의 결합 모서리 최대 3,413·3,768 J(E44), 가지 전류 공칭 약 262·255 A(결합 모서리 최대 약 262.5·257.1 A) & \textbf{조건부 통과}: 15 kJ의 23·25 \%. 저항값별 허용치는 제조사 확인 \\
P5 & V\textsubscript{A} 클램프 스위치 & Infineon FF2MR12KM1P(가지마다), 대체 Wolfspeed CAB530M12BM3 & 평균손실: 공칭 최악 481.5 W(스위칭 74 \%), 결합 모서리 최대 약 662.2·621.6 W. 켬·끔 지연 각각 1.0 µs 고정 & \textbf{통과}(선정식 환산온도): 공칭 125.1·124.9 °C, 결합 모서리 최대 147.6·142.6 °C(최소 폭, E44), 한도 150 °C. 대체 131.4·131.1 °C(공칭). 단서: 열 기준점 차, 2020년 예비 데이터시트, 0.5 µs·3.5 V 반례 151.26 °C(E45) \\
P6 & 바이패스 소프트 투입 & EPC2305(바이패스 GaN) 1개 + Stage C 일시 정지 & 최악 0.036 J·104 W, 바이패스 전류 최대 45.0 A(74사례) & \textbf{통과}: 109.6 °C(규칙 L), 전류는 소자 1개 선형 한도 0.5 × 329 A = 164.5 A 안(E46; 이전 판정은 두 병렬 소자 합계 329 A를 씀) \\
P7 & H-브리지 역방향 도통(이중 고장) & EPC2305 × 2 & 소자 최대 32.7 A, 400 V 부하 정지 3 ms & \textbf{통과}: 100.6 °C(규칙 L) \\
P8 & 링크 제동 저항 3 Ω & Vishay LTO 150H & 제동 에너지 최대 51.9 J(R2-W460 930 V H-브리지 개방, 감발 메시지 모두 손실; 약 57.5 ms의 30회 펄스, 첨두 3.3 kW, E45) & \textbf{조건부 후보}: 펄스열의 열누적·반복 주기를 제조사 자료·열 모델·시험으로 확인(100 ms 단발 곡선 × 반복 계수 52.8 J와의 비교는 여유를 보이지 않음) \\
P9 & 400 V 부하 접촉기 & TE ECK100B(1000 V) & R2-W460 470 V·52.7 A, R2-W410 420 V·70.7 A & \textbf{통과}: 100 A·800 V 1,000회 \\
C\textsubscript{A} & V\textsubscript{A} 노드 커패시터 2.4 mF & TDK B32778Z5127(600 V, 120 µF) × 20 & 리플 최대 50.6 A(공차 표본), 전압 기준 530 V & \textbf{통과}: 690 A, 2,950 cm\textsuperscript{3} \\
F2 & 이중 고장 경로전류 5 ms & 400 V 부하 정지 ≤ 3 ms(방전 저항 없음) & 경로전류 1 A 아래까지 1.83–1.98 ms & \textbf{통과}(4/4, 회로 TB31) \\
\end{longtable}
}
\srcnote{주: 모든 판정은 모의 스트레스와 데이터시트 값의 대조이며 부품 시험이 아니다. 자료: 검증 근거 E34, E37–E40, 회로 모델 TB29–TB33, 각 부품 데이터시트\cite{b34,b56,b57,b58,b59,b60,b15}.}

\Needspace{10\baselineskip}\section{판정의 경과}

부품 판정은 세 번 하고 두 번 다시 판정했다(표 \ref{t:fails}의 E34·E37–E40 행, 클램프는 표 \ref{t:clamp_hist}). 1차(E34)에서는 9개 중 4개가 불통과였다(P2 C\textsubscript{fly} 전류 밀도, P5 클램프 IGBT 29–40 kHz·334–405 °C, P6 소프트 투입 176 °C, P9 W460 접촉기; 조건부 2, 요구만 1). 2차(E37)는 후보 계열을 넓혀 P2(공진용 필름), P5(1.2 Ω + IGBT, 여유 1 °C), P9(TE ECK100B)를 통과시켰지만 P6은 불통과, 이중 고장 방전 저항 1 Ω은 2/4여서 부품으로 풀지 못해 설계를 바꿨다. 3차(E38)에서는 외부 검토로 클램프 한 가지 저항이 불성립해 가지로 나눴고, P6은 Stage C 일시 정지로 108.0 °C, F2는 400 V 부하 정지 3 ms로 4/4가 되었다(새 인터페이스 요구 두 가지). 재판정 E39는 클램프 스위치를 전형 열저항으로 판정한 구현 누락(183.3 °C 불통과)을 고쳐 CAB530M12BM3(146.5 °C)로 바꿨고, 재판정 E40에서는 규칙 L로 E39 선정이 159.7 °C 불통과가 되어 클램프를 가지 3개 × 2.4 Ω + FF2MR12KM1P로 바꾸고 최대 연속 켜짐 0.5 ms와 문턱 검출 정확도 요구를 더했다(P1·P6·P7 통과, P8은 반복 펄스 기준으로 조건부, C\textsubscript{A}는 전체 표에서 다시 고름).

함께 기록할 사항은 다음과 같다.

\begin{itemize}
\item \textbf{이 표의 계산 조건에서 나온 불통과는 설계로 풀었고, 판정 기준을 낮춰 통과시킨 것은 없다.} 후속 반례(클램프 지연 0.5 µs, 충전 기동 경부하)는 해결하지 않고 한계로 남겼다. 바뀐 것은 부품 선택(P2·P5·P9·C\textsubscript{A}), 회로 구성(클램프 가지), 보호 순서(Stage C 일시 정지, 400 V 부하 정지, 최대 연속 켜짐 시간)다.
\item \textbf{판정 규칙은 검토마다 보수 쪽으로 바뀌었다.} 공칭값 → 저항 모서리 → 최대 열저항 → 규칙 L이며, 그때마다 이전 선정을 다시 판정해 불통과도 그대로 기록했다.
\item \textbf{조건부 부품은 둘(P4 클램프 저항, P8 제동 저항), 단서가 붙은 통과는 둘(클램프 스위치, P2 C\textsubscript{fly})이다.} 확인할 내용은 표 \ref{t:parts}의 판정 칸에 적었고(C\textsubscript{fly} 용량 −4·−5 \%의 동특성은 E44에서 통과), 하드웨어 단계에서 확인한다(19장).
\end{itemize}

\section{그 밖의 요구값}

부품 후보를 정하지 않고 요구값만 정한 것은 다음과 같다. H-브리지 필터 인덕터 포화 전류 ≥ 150 A(전류 100 A 기준의 위반 이력을 지우지 않음), DAB 누설 인덕터·변압기 1차 전류는 창별로 다르다: R2-W460은 회로 관측 최대 15.2 A(이전 TB12; 최종 TB13 14.4·14.0 A), R2-W410은 평균모델 강등 대표점의 2차 환산 첨두 83.42 A → 1차 약 20.35 A(권선비 4.1, 모의 수락 한도 2차 90 A ↔ 1차 21.95 A, E46). 한 점의 관측값으로 포화 전류를 확정하지 않으며 창별 실효·첨두·전압시간적·코어 재료·턴수로 정해야 한다. 강등 중 DAB 첨두는 정상 주기 지령표 상한 1.387 × 정격점과 고장 회로 순시 수락 기준 1.5배(R2-W460 68.7 A, R2-W410 90 A)를 구분한다(TB13 R2-W460 회로 첨두 66.4·64.4 A는 지령표 63.6 A를 넘고 수락 한도 안; R2-W410의 소자 회로 첨두는 미확인), 창 조정 DAB 정격 = 정적 최대 링크전력 × 1.8, 조정형 Stage C 데드타임 1차 100·2차 40 ns(창 폭 20–30 ns), Stage C 회생 차단 ≤ 0.5–1.0 ms·≥ 4 kW/ms, Stage B 게이트 분해능 ≤ 0.625 ns.
